% Portada + abstract en la misma primera página (sin salto forzado).
\begingroup
\thispagestyle{empty}
\begin{center}
\vspace*{0.4cm}
{\LARGE\bfseries Elite Family Networks and the Reproduction of\\ Political Power in Latin America\par}
\vspace{0.75em}
{\large Evidence from Wikipedia Genealogical Data,\\ Nineteenth to Twenty-First Centuries\par}
\vspace{2em}
{\large Matías Deneken\par}
\vspace{0.35em}
{\normalsize Pontificia Universidad Católica de Chile\par}
\vspace{1.1em}
{\normalsize \today\par}
\vspace{1.2em}
{\small\itshape Working paper prepared for the Computational Text Analysis and Social Science (COMPTEXT) Conference 2026, University of Birmingham.\par}
\end{center}
\vspace{0.8cm}
\begin{minipage}[t]{0.92\textwidth}
\setlength{\parskip}{0.32em}
\noindent\textbf{Abstract.}\quad\small

How do elite family networks structure the reproduction of political power in Latin America from the nineteenth to the twenty-first century? Using genealogical data from Spanish-language Wikipedia (6,711 individuals across ten countries), we construct a multi-country kinship network linking spouses, parent–child, and partner ties (8,201 matched edges), augmented by institutional layers for party affiliation, educational trajectories, and political succession.

At the surface, we document substantial cross-national heterogeneity: closure, topology, and transnational ties vary across republics, with Pacific–Andean corridors (notably Chile–Peru and Argentina–Uruguay) structuring cross-border linkages. Beneath this variation, all networks share a common generative architecture. Across multiple complementary analyses --- stochastic block modeling, exponential random graph models, homophily decomposition, temporal cohort analysis, Gould--Fernandez brokerage decomposition, and targeted-attack simulations --- we identify a highly asymmetric core--periphery structure composed of a small set of tightly connected families (21 in total, 16 of them Chilean) with dramatically higher internal connectivity, and a set of four Medici-like robust brokers (Caicedo, Cruz, Pinto, Vicuña) that bridge core and periphery.

This architecture is dynamically reinforced over time: the share of cross-family ties declines by 70\% across five birth cohorts, from the Colonial period to the Modern era, even as total edge volume rises --- the empirical signature of oligarchic consolidation. We interpret these findings as a structured description of a "documentary elite" encoded in Wikipedia: not an exhaustive map of Latin American elites, but a large-scale archive of publicly visible kinship ties. While causal claims remain limited, the results suggest that elite reproduction in the region is simultaneously heterogeneous in form and unified by a deeply unequal network structure.

\vspace{0.55em}
\noindent\small\textit{All project materials are publicly available at \url{https://github.com/matdknu/familiaR-wiki}.}

\vspace{0.75em}
\noindent\textbf{Keywords:}\quad
elite networks; Latin America; family; political reproduction; Wikipedia; social network analysis
\end{minipage}
\par\endgroup
